\section*{Bab \ref{ch:b2:liquid-vapour-diagrams} --- Diagram Biner Cair--Uap dan Distilasi}

\begin{solution}{exo:b2:liquid-vapour-diagrams:1}
Benzena \qty{80.1}{\degreeCelsius}, toluena \qty{110.6}{\degreeCelsius}.
Cairan dengan komposisi 0.2 mulai mendidih pada sekitar
\qty{102}{\degreeCelsius}; uap pertamanya berkomposisi 0.38.
\end{solution}

\begin{solution}{exo:b2:liquid-vapour-diagrams:2}
$p = 0.5 \times 1.010 + 0.5 \times 0.388 = \qty{0.699}{bar}$;
$y = 0.505/0.699 = 0.722$.
\end{solution}

\begin{solution}{exo:b2:liquid-vapour-diagrams:3}
$n_V/n = (0.30 - 0.20)/(0.45 - 0.20) = 0.40$: \qty{4.0}{mol} uap dan
\qty{6.0}{mol} cairan. Pemeriksaan: $6.0 \times 0.20 + 4.0 \times 0.45 = 3.0 = 10
\times 0.30$.
\end{solution}

\begin{solution}{exo:b2:liquid-vapour-diagrams:4}
(a) $v = 2 + 2 - 1 = 3$: $T$, $p$, dan $x$ bebas. (b) $v = 2 + 2 - 2 = 2$;
pada tekanan tetap, 1. (c) \omterm{def:b2:liquid-vapour-diagrams:azeotrope}{Azeotrop} menambahkan hubungan $x = y$: $v = 1$;
pada tekanan tetap 0 --- \omterm{def:b2:liquid-vapour-diagrams:azeotrope}{azeotrop} mendidih pada suhu tertentu, seperti
cairan murni, walaupun komposisinya berubah terhadap tekanan.
\end{solution}

\begin{solution}{exo:b2:liquid-vapour-diagrams:5}
Pada \omterm{def:b2:liquid-vapour-diagrams:bubble-dew}{titik embun} $yp/p_1^* + (1 - y)p/p_2^* = 1$. Pada
\qty{100}{\degreeCelsius}: $0.3 \times 1.013/1.80 + 0.7 \times 1.013/0.742 = 1.124 > 1$ (terlalu dingin); pada \qty{105}{\degreeCelsius}:
$0.148 + 0.824 = 0.971 < 1$. Dengan interpolasi, $T \approx
100 + 5 \times 0.124/0.153 = \qty{104}{\degreeCelsius}$. Tetes pertama
mempunyai $x =
yp/p_1^* \approx 0.3 \times 1.013/2.0 = 0.15$.
\end{solution}

\begin{solution}{exo:b2:liquid-vapour-diagrams:6}
Tetes-tetes pertama mempunyai komposisi uap di atas cairan yang mendidih,
$y = 0.71$, pada \qty{92.1}{\degreeCelsius}. Karena benzena pergi lebih
dahulu, cairan makin miskin benzena, titik didihnya naik, dan distilatnya
juga makin miskin; tetes-tetes terakhir hampir berupa toluena murni.
Distilasi sederhana adalah satu tahap kesetimbangan, yang diulang pada
cairan yang terus berubah: setiap fraksi berupa campuran.
\end{solution}

\begin{solution}{exo:b2:liquid-vapour-diagrams:7}
Umpan 0.05, di sisi air dari \omterm{def:b2:liquid-vapour-diagrams:azeotrope}{azeotrop}: puncak menghasilkan \omterm{def:b2:liquid-vapour-diagrams:azeotrope}{azeotrop} (0.88
dalam mol, 95~\% massa), pendidih mempertahankan air. Umpan 0.95, di sisi
etanol: puncak tetap menghasilkan \omterm{def:b2:liquid-vapour-diagrams:azeotrope}{azeotrop}, titik didih terendah, dan
pendidih mempertahankan etanol murni.
\end{solution}

\begin{solution}{exo:b2:liquid-vapour-diagrams:8}
Campuran itu mendidih ketika $0.10 + p^*_{\text{air}} = 1.013$, yaitu
$p^*_{\text{air}} = \qty{0.91}{bar}$, di dekat \qty{97}{\degreeCelsius}.
Rasio massa $0.10 \times 136/(0.91 \times 18.02) = 0.83$: \qty{0.83}{g}
esens per gram air.
\end{solution}

\begin{solution}{exo:b2:liquid-vapour-diagrams:9}
Pada suhu kamar, dua lapisan L$_2$ dan L$_1$ (komposisi 0.4 terletak di
dalam celah). Ketika dipanaskan, kedua lapisan mendidih bersama pada suhu
\omterm{def:b2:liquid-vapour-diagrams:miscibility-gap}{heteroazeotrop}, yang tetap; uap pertama mempunyai komposisi \omterm{def:b2:liquid-vapour-diagrams:miscibility-gap}{heteroazeotrop}
(titik itu). Karena 0.4 terletak di sisi L$_2$ dari titik itu, lapisan
L$_1$ habis lebih dahulu; suhu kemudian naik, L$_2$ yang tersisa makin
miskin akan 1 di sepanjang \omterm{def:b2:liquid-vapour-diagrams:bubble-dew}{kurva gelembungnya} dan uap mengikuti \omterm{def:b2:liquid-vapour-diagrams:bubble-dew}{kurva
embun}, sampai titik itu mencapai \omterm{def:b2:liquid-vapour-diagrams:bubble-dew}{kurva embun}: tetes terakhir menguap.
\end{solution}

\begin{solution}{exo:b2:liquid-vapour-diagrams:10}
$N_{\min} = \ln(99 \times 99)/\ln 2.47 = 9.19/0.904 = 10.2$, dibandingkan
dengan 6.5: beralih dari kemurnian 95~\% ke 99~\% di kedua ujung menuntut
lebih dari separuh jumlah pelat sebelumnya sebagai tambahan. Setiap faktor
tambahan pada $x/(1 - x)$ menuntut jumlah pelat yang sama, dan kemurnian
diukur dari pengotor yang tersisa, yang hanya berkurang secara geometris.
\end{solution}

\begin{solution}{exo:b2:liquid-vapour-diagrams:11}
\omterm{def:b2:liquid-vapour-diagrams:binary-diagram}{Diagram isobarik} saja tidak memberikannya: diperlukan $p_1^*$ dan $p_2^*$
pada \qty{80}{\degreeCelsius}, dari persamaan Antoine (1.010 dan
\qty{0.388}{bar}). Lalu $p = p_2^* + x(p_1^* - p_2^*)$ dan
$p = 1/(y/p_1^* + (1 -
y)/p_2^*)$. Cairan stabil pada tekanan tinggi
(memampatkan uap mengembunkannya), dan pada suhu rendah pada \omterm{def:b2:liquid-vapour-diagrams:binary-diagram}{diagram
isobarik}: kedua diagram itu terbalik satu terhadap yang lain.
\end{solution}

\begin{solution}{exo:b2:liquid-vapour-diagrams:12}
Karena $y$ bertambah seiring $x$ dan $y(1 - y) > 0$, $dp/dx$ bertanda sama
dengan $y - x$: tekanan naik seiring $x$ di tempat uap lebih kaya akan 1.
Untuk \omterm{def:b2:chemical-potential:ideal-mixture}{campuran ideal}, $y > x$ untuk setiap $0 < x < 1$
(\cref{prop:b2:liquid-vapour-diagrams:richer}): $p$ monoton tegas, tidak
mempunyai ekstremum, dan tidak ada \omterm{def:b2:liquid-vapour-diagrams:azeotrope}{azeotrop}.
\end{solution}

\begin{solution}{pb:b2:liquid-vapour-diagrams:1}
\textbf{1.} $T = B/(A - \log_{10}1.013) - C$: benzena $1203.835/4.01253 + 53.226 =
\qty{353.2}{K}$, \qty{80.1}{\degreeCelsius}; toluena $1343.943/4.07266 + 53.773 =
\qty{383.8}{K}$, \qty{110.6}{\degreeCelsius}.
\textbf{2.} Pada \qty{363.15}{K}: $p_1^* = \qty{1.361}{bar}$,
$p_2^* = \qty{0.542}{bar}$.
\textbf{3.} $x = (1.013 - 0.542)/(1.361 - 0.542) = 0.575$; $y = 0.575 \times
1.361/1.013 = 0.773$.
\textbf{4.} $x = (1.013 - 0.742)/(1.800 - 0.742) = 0.256$; $y = 0.256 \times
1.800/1.013 = 0.455$.
\textbf{5.} Titik-titik gelembung (0, 110.6), (0.256, 100), (0.575, 90),
(1, 80.1); titik-titik embun (0, 110.6), (0.455, 100), (0.773, 90),
(1, 80.1) --- lensa pada gambar.
\textbf{6.} $\frac12(p_1^* + p_2^*) = 1.013$: pada \qty{90}{\degreeCelsius}
0.952, pada \qty{95}{\degreeCelsius} 1.102; dengan interpolasi,
$90 + 5 \times 0.061/0.150 \approx
\qty{92}{\degreeCelsius}$
(tepatnya \qty{92.1}{\degreeCelsius}).
\textbf{7.} $v = 2$, menjadi 1 pada tekanan tetap: setiap suhu menetapkan
$x$ dan $y$; itulah sebabnya diagram itu terdiri atas dua kurva.
\textbf{8.} $x = (1.013 - 0.636)/(1.569 - 0.636) = 0.404$; $y = 0.404 \times
1.569/1.013 = 0.626$.
\textbf{9.} $n_V = 100 \times (0.500 - 0.404)/(0.626 - 0.404) = \qty{43}{mol}$,
$n_L = \qty{57}{mol}$.
\textbf{10.} Uap $43 \times 0.626 = \qty{26.9}{mol}$; cairan $57 \times 0.404
= \qty{23.0}{mol}$; total \qty{49.9}{mol} $\approx 50$.
\textbf{11.} $26.9/50 = 54~\%$ benzena, dalam uap yang hanya 63~\%.
\textbf{12.} Pada \qty{100}{\degreeCelsius}, uap dalam kesetimbangan
mempunyai $y =
0.455 < 0.5$: titik umpan terletak di atas \omterm{def:b2:liquid-vapour-diagrams:bubble-dew}{kurva embun},
umpan seluruhnya berupa uap dan tidak ada yang terpisah.
\textbf{13.} Pada \omterm{def:b2:liquid-vapour-diagrams:bubble-dew}{titik embunnya}, \qty{98.8}{\degreeCelsius} (dengan metode
pada \cref{exo:b2:liquid-vapour-diagrams:5}).
\textbf{14.} $\alpha = p_1^*/p_2^*$; pada \qty{80.1}{\degreeCelsius},
$1.013/0.390 =
2.60$.
\textbf{15.} Pada \qty{110.6}{\degreeCelsius}, $2.377/1.013 = 2.35$.
\textbf{16.} $\sqrt{2.60 \times 2.35} = 2.47$.
\textbf{17.} $y = xp_1^*/p$ dan $1 - y = (1 - x)p_2^*/p$; bagilah.
\textbf{18.} Tidak ada yang diambil: di antara dua pelat, uap yang naik dan
cairan yang turun membawa jumlah materi dan jumlah setiap komponen yang
sama, $V y_k = L x_{k+1}$ dengan $V = L$, sehingga $x_{k+1} = y_k$.
\textbf{19.} Dengan $r = x/(1 - x)$: $r(y_k) = \alpha\,r(x_k)$ dan
$x_{k+1} = y_k$, sehingga $r(x_{k+1}) = \alpha\,r(x_k)$ dan
$r(x_D) = \alpha^N r(x_B)$; $N = \ln[r(x_D)/
r(x_B)]/\ln\alpha$.
\textbf{20.} $N_{\min} = \ln(19 \times 19)/\ln 2.47 = 5.889/0.904 = 6.5$.
\textbf{21.} Kolom harus menghasilkan produk, sehingga dijalankan pada \omterm{def:b2:liquid-vapour-diagrams:fractional-distillation}{rasio
refluks} berhingga, yang memerlukan lebih banyak pelat; selain itu, pelat
nyata tidak mencapai kesetimbangan (efisiensinya kurang dari satu).
\textbf{22.} Tujuh anak tangga, yang terakhir berakhir pada 0.034, di bawah
0.05: 6.5 pelat dalam hitungan kontinu, tujuh pelat utuh (pendidih ulang
dihitung sebagai satu).
\textbf{23.} $x = (95/46.07)/(95/46.07 + 5/18.02) = 0.881$.
\textbf{24.} Di puncak, paling banyak \omterm{def:b2:liquid-vapour-diagrams:azeotrope}{azeotrop} (0.88 dalam mol); di dasar,
air. Lima puluh pelat tidak melintasi \omterm{def:b2:liquid-vapour-diagrams:azeotrope}{azeotrop}.
\textbf{25.} Pada refluks total, $\boldsymbol{N_{\min} \approx 6.5}$ \omterm{def:b2:liquid-vapour-diagrams:fractional-distillation}{pelat
teoretis}.
\end{solution}
